\section{Results and intended use}
\label{sec:overview}

The useful extension is a way to produce compatible mathematical
observations from a system of numerical dependencies. It applies to
simultaneous equations, approximate inverses, interacting analytic
constructions, and finite models whose size increases as more information
is requested. Its input is finite algebra and explicit error estimates;
its output is a primitive recursive sampler with a proved modulus.

There are two central statements. First, suppose a rational implementation
\(E_p\) maps a certified box into itself and satisfies, componentwise,
\[
 |E_p(x)-E_r(y)|\le A|x-y|+(2^{-p}+2^{-r})w,
 \qquad Aw\le(1-1/s)w,
 \tag{1}\label{eq:overview-action}
\]
where \(A\ge0\), \(w_i>0\), and \(s\ge1\) is an integer. The weighted
distance is \(d_w(x,y)=\max_i |x_i-y_i|/w_i\). A finite iteration with an
explicit precision schedule constructs \(Q(k)\) satisfying
\[
 d_w(Q(k),Q(l))\le2^{-k}+2^{-l}.
 \tag{2}\label{eq:overview-name}
\]
Moreover, its residual tends to zero at a supplied rate, and any other
small-residual computation agrees quantitatively with this one. There
is no assumption that a completed real fixed point has already been
constructed. The maps \(E_p\) may be discontinuous rounding procedures;
the mixed estimate in (1) is the property used in the proof.

Second, suppose finite-level updates \(F_n\) and restrictions \(P_n\)
satisfy a commutation estimate
\[
 d_n(F_nP_n z,P_nF_{n+1}z)\le\eta_n,
\]
and the actual returned points \(y_n\) have residual at most \(r_n\).
If \(F_n\) contracts by \(q_n<1\), then
\[
 d_n(y_n,P_ny_{n+1})\le
 \frac{r_n+\eta_n+r_{n+1}}{1-q_n}
 \tag{3}\label{eq:overview-refinement}
\]
for nonexpansive restrictions. An explicit summable bound on the right
produces a coherent family of observations. This supplies the missing
quantitative connection between solving a finite model and obtaining one
object across successive models.

\subsection{A concrete infinite-dimensional endpoint}

For \(i\in\N\), consider
\[
 x_i=\frac1{4(i+1)}+\frac{x_i^2}{8}+\frac{x_{i+1}}8,
 \qquad 0\le x_i\le1.                         \tag{4}\label{eq:overview-lattice}
\]
The report constructs a single primitive recursive rational function
\(X(i,k)\). To answer one query it uses only
\(i+k+3\) coordinates, \(i+k+7\) rounded iterations, and
\(i+k+11\) binary fractional digits. It proves
\[
 \begin{split}
 |X(i,k)-X(i,l)|&\le2^{-\min(k,l)-2},\\
 \left|X(i,k)-\frac1{4(i+1)}-
 \frac{X(i,k)^2}{8}-\frac{X(i+1,k)}8\right|&\le2^{-k-4}.
 \end{split}                                                  \tag{5}
\]
These are uniform arithmetic assertions, not a finite table of test
values. They express a unique solution to (4) within the specified box.
Uniqueness is not asserted among all bounded sequences.

\subsection{What is proved, and at what level}

\begin{longtable}{@{}>{\raggedright\arraybackslash}p{0.27\textwidth}>{\raggedright\arraybackslash}p{0.41\textwidth}>{\raggedright\arraybackslash}p{0.24\textwidth}@{}}
\toprule
Result & Mathematical endpoint & Evidence in this packet\\
\midrule
\endhead
Coupled approximation & Explicit residual bound and regular sampler from
(1) & Full human proof; arithmetic core compiled in Lean\\
Matrix certificates & Positive weights, componentwise supersolutions,
finite Neumann bounds & Full finite proofs and exact reference checks\\
Polynomial systems & Root names and parameter stability from a finite
box certificate & Full human proof and explicit coefficient algorithms\\
Refinement towers & Coherence from residuals, commutation defects and
contraction gaps & Full finite proof and formalization definitions\\
Infinite nonlinear system & PR coordinate sampler, equations and uniqueness
on \([0,1]^{\N}\) & Full human proof and exact reference computation\\
Object-theory export & \(\Gamma\vdash[\texttt{PRAAx}]\varphi\) for the
coded endpoints & Detailed invariant plan and an uncompiled adapter\\
\bottomrule
\end{longtable}

The human arguments identify the primitive recursive definitions and
quantifier-free induction formulas required for iPRA. The compiled
Lean arithmetic lemmas check important finite identities and bounds.
They do not, by themselves, certify the entire syntactic translation into
\texttt{Prf PRAAx}. This separation is part of the deliverable's exact
claim, rather than a condition left for the reader to discover.

The underlying contraction and verified-numerics ideas are established
mathematics; compare Rump's finite-system treatment and the radii-polynomial
method \cite{rump1994,radii}. The contribution here is the particular
finitary organization, the approximate-action and cross-level interfaces,
the explicit schedules, and their integration with the existing iPRA
programme. No claim of priority for the general contraction principle is
made.
